% --- 假设与符号 ---
\section{假设与符号}\label{sec:assumptions}

\subsection{数据口径与单位}
为保证模型求解与结果可比，采用如下统一口径：
\begin{itemize}
  \item 面积单位：亩（或公顷，本文以“亩”为例）；产量单位：kg；价格单位：元/kg；成本单位：元/亩；销量单位：kg。
  \item 时标：年 $t\in\mathcal{T}=\{2024,\ldots,2030\}$，季 $s\in\mathcal{S}$（如春/夏/秋/冬，或早/晚季，依数据口径）。
  \item 地块 $i\in\mathcal{I}$，其中大棚地块子集 $\mathcal{G}\subseteq\mathcal{I}$；作物 $k\in\mathcal{K}$。
  \item 以 2023 年为基期，成本、价格、单产等参数若无特殊说明均以名义值计。
\end{itemize}

\subsection{基本假设}
\begin{enumerate}
  \item 地块可用性：给定 $A_{is}$ 表示地块 $i$ 在季 $s$ 的可用面积；季节-地块的不可用状态用 $A_{is}=0$ 表示。
  \item 作物适宜性：$\eta_{isk}\in\{0,1\}$ 指示作物 $k$ 在地块 $i$ 的季 $s$ 是否可种（已将大棚专属限制并入 $\eta_{isk}$ 或集合 $\mathcal{K}_{\mathrm{green}}$）。
  \item 产量线性化：单位面积产量 $y_{tsk}$ 给定，季节与年份可变；总产量等于 $y_{tsk}\sum_i x_{tsik}$。
  \item 价格/成本：确定性场景下单价 $p_{tsk}$、单位面积成本 $c_{tsk}$ 为常数；销量上限 $u_{tsk}$ 给定。
  \item 市场可售：若无特殊折价机制，先以“基准价—销量上限”框架求解；超产折价若存在在评价节（或扩展模型）中再引入。
  \item 劳作资源与机械约束：如未提供详表，先不显式建模，后续可按“年/季最大可耕面积或工时”扩展。
\end{enumerate}

\subsection{农艺规则与可行域补充}
\begin{itemize}
  \item 同季不重叠：同一地块同一季仅允许一种作物。
  \item 重茬与轮作：同一地块禁止相邻季对同一作物的连续种植；豆科等作物遵循最小轮作间隔 $L_k$（如 $L_k\in\{2,3\}$ 季）。
  \item 大棚限制：若 $i\notin\mathcal{G}$，则 $\mathcal{K}_{\mathrm{green}}$ 中作物在该地块不可种；已体现在 $\eta_{isk}$ 中。
  \item 面积上限：$0\le x_{tsik}\le A_{is}$ 且受适宜性 $\eta_{isk}$ 与二元选择变量联动。
\end{itemize}

\subsection{价格与销量机制}
\begin{itemize}
  \item 确定性：$0\le q_{tsk}\le u_{tsk}$，收入为 $p_{tsk}\,q_{tsk}$；若 $q_{tsk}$ 受产量约束，则 $q_{tsk}\le \sum_i y_{tsk}\,x_{tsik}$。
  \item 可选折价（扩展）：当 $\sum_i y_{tsk}\,x_{tsik}>u_{tsk}$，超出部分按折减系数 $\phi_{tsk}\in(0,1)$ 定价：$p_{tsk}^{\mathrm{eff}}=p_{tsk}$（不超时），超出部分单价 $\phi_{tsk}p_{tsk}$。
\end{itemize}

\subsection{符号表（扩展）}\label{sec:assumptions-symbols}
为便于查阅，这里在“极简符号表”的基础上给出更完整的记号分组与说明（单位/口径已与本章前述约定一致）。

\paragraph{集合与索引}
\begin{table}[H]\centering
\begin{tabular}{@{}ll@{}}\toprule
\(\mathcal{T}\) & 年集合（2024--2030） \\
\(\mathcal{S}\) & 季节集合（按数据口径定义顺序，如春\(\to\)夏\(\to\)秋\(\to\)冬） \\
\(\mathcal{I}\) & 地块集合 \\
\(\mathcal{K}\) & 作物集合 \\
\(\mathcal{G}\subseteq\mathcal{I}\) & 大棚地块子集（如适用） \\
\(\mathcal{K}_{\mathrm{green}}\subseteq\mathcal{K}\) & 仅允许在大棚种植的作物集合（如适用） \\
\(\mathcal{K}_{\mathrm{legume}}\subseteq\mathcal{K}\) & 需遵循最小轮作间隔的作物集合（如豆科） \\
\(\mathcal{W}\) & 情景集合（用于滚动/仿真/鲁棒评估） \\
\bottomrule\end{tabular}
\end{table}
\noindent 说明：跨年滚动时季节序按数据给定口径循环推进。

\paragraph{决策变量}
\begin{table}[H]\centering
\begin{tabular}{@{}ll@{}}\toprule
\(x_{t s i k}\ge 0\) & 年\(t\)、季\(s\)、地块\(i\)上作物\(k\)的种植面积（亩） \\
\(z_{t s i k}\in\{0,1\}\) & 是否在\((t,s,i)\)上选择作物\(k\)（控制同季不重叠/重茬） \\
\(q_{t s k}\ge 0\) & 年\(t\)、季\(s\)作物\(k\)的销售量（kg） \\
\(r_t\) & 年\(t\)净收益（统计用） \\
\bottomrule\end{tabular}
\end{table}

\paragraph{参数}
\begin{table}[H]\centering
\begin{tabular}{@{}ll@{}}\toprule
\(A_{i s}\) & 地块\(i\)在季\(s\)可用面积（亩） \\
\(\eta_{i s k}\in\{0,1\}\) & 作物\(k\)在地块\(i\)季\(s\)的适宜性指示 \\
\(y_{t s k}\) & 单产（kg/亩） \\
\(p_{t s k}\) & 单价（元/kg） \\
\(c_{t s k}\) & 单位面积成本（元/亩，含种植/管理等） \\
\(u_{t s k}\) & 销量上限（kg） \\
\(L_k\) & 作物\(k\)的最小轮作间隔（以季计） \\
\(\phi_{t s k}\in(0,1)\) & 超产折价系数（可选，若启用折价机制） \\
\(\Gamma\) & 预算不确定集保守度（鲁棒优化中使用） \\
\(\alpha\in(0,1)\) & CVaR 置信水平（如 0.95，用于尾部风险度量） \\
\bottomrule\end{tabular}
\end{table}

\noindent 注：若题面未启用折价与鲁棒/动态机制，可将 \(\phi_{t s k}\)、\(\Gamma\)、\(\alpha\) 视为未激活参数；数据来源见题目附件与数据预处理节。